\section{回顾：从线性分类器到优化}

讲者 Ehsan Adeli 首先回顾了前三节课的核心内容，为本节课的神经网络和反向传播奠定基础。

\begin{figure}[H]
\centering
\includegraphics[width=0.95\textwidth]{slides-images/slide-001.jpg}
\caption{课程回顾：从评分函数到损失函数再到优化的完整流程\protect\footnotemark}
\end{figure}
\footnotetext{Slides 第2页。}

\subsection{评分函数与损失函数}

我们已经建立了完整的学习框架：

\begin{enumerate}
    \item \textbf{评分函数}：线性模型 $f(\mathbf{x}, W) = W\mathbf{x}$，将输入映射到类别分数
    \item \textbf{损失函数}：衡量预测的好坏，包括数据损失和正则化项
    \item \textbf{优化}：通过梯度下降找到最优权重 $W$
\end{enumerate}

\begin{figure}[H]
\centering
\includegraphics[width=0.95\textwidth]{slides-images/slide-003.jpg}
\caption{学习流程图：输入 $\mathbf{x}$ 经过评分函数得到分数 $s$，与标签 $y$ 一起计算损失 $L$\protect\footnotemark}
\end{figure}
\footnotetext{Slides 第4页。}

\begin{knowledgebox}{Hinge Loss vs Softmax Loss}
除了上一节课详细讨论的 Softmax（交叉熵）损失外，另一种常见的分类损失是 \textbf{Hinge Loss}（SVM 损失）：
$$L_i = \sum_{j \neq y_i} \max(0, s_j - s_{y_i} + 1)$$
它鼓励正确类别的分数比所有错误类别的分数至少高出一个 margin（此处为 1）。与 Softmax 不同，Hinge Loss 不将分数转换为概率。
\end{knowledgebox}

\subsection{优化回顾}

\begin{figure}[H]
\centering
\includegraphics[width=0.95\textwidth]{slides-images/slide-007.jpg}
\caption{梯度下降：在损失景观中沿负梯度方向逐步移动到最优点\protect\footnotemark}
\end{figure}
\footnotetext{Slides 第8页。}

梯度下降的更新规则为 $W \leftarrow W - \alpha \nabla_W L$。为了降低计算成本，我们使用 \textbf{SGD}（随机梯度下降），每次只用一个 mini-batch（如 32、64、128 或 256 个样本）来估计梯度。我们还讨论了 SGD+Momentum、RMSProp 和 Adam 等高级优化器。

\begin{figure}[H]
\centering
\includegraphics[width=0.95\textwidth]{slides-images/slide-011.jpg}
\caption{SGD 及其改进：Momentum、RMSProp 和 Adam 优化器的对比\protect\footnotemark}
\end{figure}
\footnotetext{Slides 第12页。}

\subsection{本章小结}

前几节课建立了``评分函数 $\rightarrow$ 损失函数 $\rightarrow$ 梯度下降''的完整框架，但仅限于线性分类器。本节课将扩展到非线性的神经网络，并介绍如何高效计算复杂网络中的梯度——这就是反向传播算法。

%% ========== 神经网络 ==========

